%----------
%   WARNING
%----------

% This Guide contains Library recommendations based mainly on APA and IEEE styles, but you must always follow the guidelines of your TFG Tutor and the TFG regulations for your degree.

% THIS TEMPLATE IS BASED ON THE APA STYLE 


%----------
% DOCUMENT SETTINGS
%----------

\documentclass[12pt]{report} % font: 12pt

% margins: 2.5 cm top and bottom; 3 cm left and right
\usepackage[
a4paper,
vmargin=2.5cm,
hmargin=3cm
]{geometry}

% Paragraph Spacing and Line Spacing: Narrow (6 pt / 1.15 spacing) or Moderate (6 pt / 1.5 spacing)
\renewcommand{\baselinestretch}{1.15}
\parskip=6pt

% Color settings for cover and code listings 
\usepackage[table]{xcolor}
\definecolor{azulUC3M}{RGB}{0,0,102}
\definecolor{gray97}{gray}{.97}
\definecolor{gray75}{gray}{.75}
\definecolor{gray45}{gray}{.45}

% In the template we include the file OUTPUT.XMPDATA. You can download that file and include the metadata that will be incorporated into the PDF file when you compile the memoria.tex file. Then upload it back to your project. 
% Commented out due to upstream error
% \use package[a-1b]{pdf}

% LINKS
\usepackage{hyperref}
\hypersetup{colorlinks=true,
	linkcolor=black, % links to parts of the document (e.g. index) in black
	urlcolor=blue} % links to resources outside the document in blue

% MATH EXPRESSIONS
\usepackage{amsmath,amssymb,amsfonts,amsthm}

% Character encoding
\usepackage{txfonts} 
\usepackage[T1]{fontenc}
\usepackage[utf8]{inputenc}

% English settings
\usepackage[spanish]{babel} 
\usepackage[babel, spanish=spanish]{csquotes}
\AtBeginEnvironment{quote}{\small}

% Footer settings
\usepackage{fancyhdr}
\pagestyle{fancy}
\fancyhf{}
\renewcommand{\headrulewidth}{0pt}
\rfoot{\thepage}
\fancypagestyle{plain}{\pagestyle{fancy}}

% DESIGN OF THE TITLES of the parts of the work (chapters and epigraphs or sub-chapters)
\usepackage{titlesec}
\usepackage{titletoc}
\titleformat{\chapter}[block]
{\large\bfseries\filcenter}
{\thechapter.}
{5pt}
{\MakeUppercase}
{}
\titlespacing{\chapter}{0pt}{0pt}{*3}
\titlecontents{chapter}
[0pt]                                               
{}
{\contentsmargin{0pt}\thecontentslabel.\enspace\uppercase}
{\contentsmargin{0pt}\uppercase}                        
{\titlerule*[.7pc]{.}\contentspage}                 

\titleformat{\section}
{\bfseries}
{\thesection.}
{5pt}
{}
\titlecontents{section}
[5pt]                                               
{}
{\contentsmargin{0pt}\thecontentslabel.\enspace}
{\contentsmargin{0pt}}
{\titlerule*[.7pc]{.}\contentspage}

\titleformat{\subsection}
{\normalsize\bfseries}
{\thesubsection.}
{5pt}
{}
\titlecontents{subsection}
[10pt]                                               
{}
{\contentsmargin{0pt}                          
	\thecontentslabel.\enspace}
{\contentsmargin{0pt}}                        
{\titlerule*[.7pc]{.}\contentspage}  


% Tables and figures settings
\usepackage{multirow} % combine cells 
\usepackage{caption} % customize the title of tables and figures
\usepackage{floatrow} % we use this package and its \ ttabbox and \ ffigbox macros to align the table and figure names according to the defined style.
\usepackage{array} % with this package we can define in the following line a new type of column for tables: custom width and centered content
\newcolumntype{P}[1]{>{\centering\arraybackslash}p{#1}}
\DeclareCaptionFormat{upper}{#1#2\uppercase{#3}\par}
\usepackage{graphicx}
\graphicspath{{../imagenes/}} % images folder

% Table layout for social sciences and humanities
\captionsetup*[table]{
	justification=raggedright,
	labelsep=newline,
	labelfont=small,
	singlelinecheck=false,
	labelfont=bf,
	font=small,
	textfont=it
}

% Figure layout for social sciences and humanities
\captionsetup[figure]{
	%name=Figura,
	singlelinecheck=off,
	labelsep=newline,
	font=small,
	labelfont=bf,
	textfont=it
}
\floatsetup[figure]{
    style=plaintop,
    heightadjust=caption,
    footposition=bottom,
    font=small
}

% Figures and tables footnote layout 
\captionsetup*[floatfoot]{
    footfont={small, up}
}

% FOOTNOTES
\usepackage{chngcntr} % continuous numbering of footnotes
\counterwithout{footnote}{chapter}

% CODE LISTINGS 
% support and styling for listings. More information in  https://es.wikibooks.org/wiki/Manual_de_LaTeX/Listados_de_código/Listados_con_listings
\usepackage{listings}

% Custom listing
\lstdefinestyle{estilo}{ frame=Ltb,
	framerule=0pt,
	aboveskip=0.5cm,
	framextopmargin=3pt,
	framexbottommargin=3pt,
	framexleftmargin=0.4cm,
	framesep=0pt,
	rulesep=.4pt,
	backgroundcolor=\color{gray97},
	rulesepcolor=\color{black},
	%
	basicstyle=\ttfamily\footnotesize,
	keywordstyle=\bfseries,
	stringstyle=\ttfamily,
	showstringspaces = false,
	commentstyle=\color{gray45},     
	%
	numbers=left,
	numbersep=15pt,
	numberstyle=\tiny,
	numberfirstline = false,
	breaklines=true,
	xleftmargin=\parindent
}

\captionsetup*[lstlisting]{font=small, labelsep=period}
 
\lstset{style=estilo}
\renewcommand{\lstlistingname}{\uppercase{Código}}


% REFERENCES 

% APA bibliography setup
\usepackage[style=apa, backend=biber, natbib=true, hyperref=true, uniquelist=false, sortcites]{biblatex}

\addbibresource{referencias.bib} % The references.bib file in which the bibliography used should be

% Caption package, for use of subfigures.
\usepackage{subfig}


%-------------
%	DOCUMENT
%-------------

\begin{document}

\pagenumbering{roman} % Roman numerals are used in the numbering of the pages preceding the body of the work.
	
%----------
%	COVER
%----------	
\begin{titlepage}
	\begin{sffamily}
  \begin{figure}%
    \raggedleft
    \subfloat{{\includegraphics[width=3cm]{UNED.jpg} }}%
    \hspace*{\fill}
    \subfloat{{\includegraphics[width=5cm]{Scalefast.jpg} }}%
\end{figure}
	\begin{center}
		\vspace{2.5cm}
		\begin{Large}
			Master en Ciberseguridad\\			
			 2020-2021\\
			\vspace{2cm}		
			\textsl{Trabajo de final de Master}
			\bigskip
			
		\end{Large}
		 	{\Huge ``Definicion, creacion e implementaicon del SecDevOps''}\\
		 	\vspace*{0.5cm}
	 		\rule{10.5cm}{0.1mm}\\
			\vspace*{0.9cm}
			{\LARGE Sergio Rosello Morell}\\ 
			\vspace*{1cm}
		\begin{Large}
			David Aracil Cofrade\\
			Rafael Pastor Vargas\\
			Madrid, a \today \\
		\end{Large}
	\end{center}
	\vfill
	\color{black}
	\fbox{
	\begin{minipage}{\linewidth}
    	\textbf{AVOID PLAGIARISM}\\
    	\footnotesize{The University uses the \textbf{Turnitin Feedback Studio} for the delivery of student work. This program compares the originality of the work delivered by each student with millions of electronic resources and detects those parts of the text that are copied and pasted. Plagiarizing in a TFM is considered a  \textbf{\underline{Serious Misconduct}}, and may result in permanent expulsion from the University.}\end{minipage}}

	% IF OUR WORK IS TO BE PUBLISHED UNDER A CREATIVE COMMONS LICENSE, INCLUDE THESE LINES. IS THE RECOMMENDED OPTION.
	\noindent\includegraphics[width=4.2cm]{creativecommons.png}\\ % Creative Commons Logo
    \footnotesize{This work is licensed under Creative Commons \textbf{Attribution – Non Commercial – Non Derivatives}}
	
	\end{sffamily}
\end{titlepage}

\newpage % blank page
\thispagestyle{empty}
\mbox{}




%----------
%	ABSTRACT AND KEYWORDS 
%----------	
\renewcommand\abstractname{\large\bfseries\filcenter\uppercase{Sumario}}
\begin{abstract}
\thispagestyle{plain}
\setcounter{page}{3}
	
Inicio del movimiento DevOps, el 2010, con el desarrollo agil\\
Adopcion de la cultura agil por silicon valley y resto de munto\\
Implementacion de seguridad dentro del desarrollo agil\\
Uso de Pipelines, para desarrollar SW con la seguridad en mente\\
Caso especifico de la implementacion en Scalefast\\
\vfill
\end{abstract}
	
\newpage % Blank page
\renewcommand\abstractname{\large\bfseries\filcenter\uppercase{Abstract}}
\begin{abstract}
\thispagestyle{plain}
\setcounter{page}{5}

Start of Agile SW development, and DevOps culture\\
Adoption of culture by Silicon Valley and subsequent popularization\\
Security inside agile development?\\
Use of Pipelines in security-focused Agile development\\
Specific case with Scalefast\\

\textbf{Palabras clave:} % add the keywords
	
\vfill
\end{abstract}
\newpage % Blank page
\thispagestyle{empty}
\mbox{}


%----------
%	Dedication
%----------	
\chapter*{Agradecimientos}

\setcounter{page}{7}
	
Agradecer principalmente este trabajo a:

\begin{itemize}
  \item{Familia}
  \item{Tutores}
  \item{UNED}
  \item{Scalefast}
\end{itemize}

		
	\vfill
	
	\newpage % blank page
	\thispagestyle{empty}
	\mbox{}
	




%----------
%	TOC
%----------	

%--
% TOC
%-
\tableofcontents
\thispagestyle{fancy}

\newpage % blank page
\thispagestyle{empty}
\mbox{}




%--
% List of figures. If they are not included, comment the following lines
%-
\listoffigures
\thispagestyle{fancy}

\newpage % blank page
\thispagestyle{empty}
\mbox{}




%--
% List of tables. If they are not included, comment the following lines
%-
\listoftables
\thispagestyle{fancy}

\newpage % blankpage
\thispagestyle{empty}
\mbox{}





%----------
%	THESIS
%----------	
\clearpage
\pagenumbering{arabic} % numbering with Arabic numerals for the rest of the document.	

\chapter{Estado del arte}

\section{Evolución del desarrollo de software, desde su concepción hasta la actualidad}

El software solo es posible gracias al hardware en el que se ejecuta.
Naturalmente, este campo ha adoptado las costumbres y características del mundo del Hardware.

Inicialmente, en 1950, actualizar el software de un ordenador suponía un gasto tremendo, ya que el software que corría el ordenador estaba estrictamente acoplado al hardware que hacia posible dicha ejecución.

Al no disponer de una conexión a Internet y requerir la presencia de ingenieros enviados por parte de la compañía que fabrica el propio ordenador, las empresas que poseían un ordenador, no esperaban tener que actualizar su software regularmente.

A medida que se avanzaba en la fabricación de ordenadores personales y se estandarizaba Internet, los requisitos para poder actualizar el software del ordenador empezaban a ofrecer menos resistencia.

A medida que evoluciona Internet, también lo hacen las aplicaciones que se alojan en el, pasamos de Internet 1.0, a  Internet 2.0, en el que las paginas web dejan de ser portales de visita, como revistas, y se empieza a poder diferenciar clientes, y adaptar el funcionamiento de la pagina al cliente especifico.

Estos nuevos requisitos, incorporan nueva complejidad a la ya creciente complejidad del software, como por ejemplo, cualquier pagina que tenga que poder diferenciar a sus clientes, por lógica de negocio, debe añadir a su pila de componentes, una base de datos.

Debido a esta creciente complejidad, las empresas cuyo modelo de negocio estaba basado en el software, ofrecían actualizaciones muy poco regularmente ya que el proceso de construcción y compilación de dicho software era tedioso.
Como añadido, a mayor complejidad de aplicación, mayor era el tiempo que se tenia que designar a la compilación de las distintas piezas que componían el software. (Base de datos, Portal frontal, Sistema de lógica interna, entre otras.)

Debido a la complejidad de las aplicaciones generadas, el proceso de Despliegue, ha sido, y sigue siendo, uno de los puntos de estrés de muchas empresas.

La realidad, es que desde 1957, en el que se pretendía desarrollar iterativamente e incrementalmente \cite{IID}, se ha estado pensando colectivamente en como mejorar el desarrollo de software, desde su fase de concepción, a su fase de despliegue, pero hasta 2001, no se escribió el manifesto del desarrollo ágil \cite{agile}, en el que se destacan los pilares fundamentales sobre los que se ha construido la filosofía DevOps \cite{CD-TF}, posteriormente extendida a la filosofía DevSecOps.

Esto último es lo que se va a proceder a analizar en este trabajo.


\subsection{Historia de los despliegues}

En esta sección analizaremos las distintas formas en las que se ha desarrollado software comúnmente, hasta hoy en día, haciendo énfasis en el desarrollo ágil, y el uso de Pipelines para garantizar cierta calidad del software antes de desplegar.

\subsubsection{1950}

En esta década, se empieza a pensar y desarrollar una mejor forma de construir software. \par
El proyecto más memorable en el que se hace uso de técnicas modernas de desarrollo de software es "Project Mercury". \par
Este proyecto, centrado en poner al hombre a la luna, requería de una capacidad de construcción de software ejemplar.
Gerald M. Weinberg fue el arquitecto del proyecto.
Decidieron que ``el desarrollo en cascada aplicado a un proyecto grande era bastante estúpido, o al menos, ignorante a la realidad'' \cite{GW-PM} y desarrollaron software con una metodología similar a la estipulada por XP.

\subsubsection{1968}

La primera aparición formal de desarrollo iterativo aparece en un documento interno de IBM en el que se describe que primero se define formalmente la funcionalidad, posteriormente se establecen unos tests, y se empieza a desarrollar, a medida que avanza el proceso, se crean nuevos tests, mas específicos, al final, el sistema, se convierte en la aplicación. \cite{ID-FB}

\subsubsection{1970}

Se documenta por primera vez el desarrollo en cascada en el articulo que escribe el Dr. Winston W. Royce. \par
Como curiosidad, Royce aconseja realizar el proceso de desarrollo compuesto por las fases de análisis de requisitos, construcción del software, testeo y despliegue dos veces, pero a medida que se ha extendido esta metodología, se ha popularizado siguiendo únicamente uno de estos flujos, quedando en el desarrollo en cascada que se conoce hoy en día. \cite{royce1970} \par
Actualmente, una de las metodologías de desarrollo más populares, debido a su simplicidad, y sentido. \par

\subsubsection{1975}

Vic Basil y Joe Turner escriben un articulo explicando como desarrollaron una serie de compiladores extendidos para una familia de lenguajes de programación específicos a dominio.
En este articulo, se describe claramente el desarrollo iterativo incremental (IID). \cite{6312870}

Durante los años 70, el desarrollo en cascada cobra una popularidad abismal, tanta, que los lideres de equipo encargados de proyectos tan importantes como la creación del software de la lanzadera espacial de la NASA, sentían la obligación de justificar por que no habían usado el método de de desarrollo en cascada.
La justificación mas común para no usar el desarrollo en cascada, era que los requisitos del proyecto cambiaban continuamente, y no podían ajustar el desarrollo en cascada a estos cambios.

\subsubsection{1980}

Durante los años 80, muchos desarrolladores prominentes publican artículos sobre como el desarrollo en cascada no es una buena forma de construir un proyecto medianamente grande de software.
Parafraseando el articulo titulado ``Un proceso de diseño racional: Como y por que falsearlo'' \cite{Parnas1986}
\begin{itemize}
  \item{Un usuario raramente sabe exactamente todo lo que quiere y no puede expresar todo lo que sabe}
  \item{Existen muchos detalles de implementación que no se pueden anticipar, aunque tengamos todos los requisitos claros}
  \item{Aunque sepamos todos los detalles, como humanos, no podemos procesar tanta complejidad}
  \item{Aunque pudiésemos procesar toda esta complejidad, existen fuerzas exteriores, que hacen que cambien los requisitos o invaliden decisiones previas}
\end{itemize}
Estas razones encabezan la lista de motivos por los que la tendencia a la hora de desarrollar software debía de virar hacia el desarrollo incremental e iterativo.

\subsubsection{1990 - Presente}

A medida que el desarrollo en cascada demostraba que no era lo suficientemente flexible para llevar a cabo proyectos medianamente grandes, cantidad de nuevos artículos aparecían con nuevas y mejoradas técnicas de desarrollo iterativo e incremental. \par
Muchas de estas siguen aplicándose hoy en día, como XP (Extreme Programming), Dynamic Systems Development Method (DSDM), Scrum y FDD (Feature Driven Development), entre otras. \par
El 1 de febrero de 2001, se celebra una reunion entre varios expertos en procesos, entre ellos, los promotores de XP, Scrum, FDD. 
En ese momento se forma la alianza ágil (agile alliance), cuya función es promover métodos de desarrollo de software iterativo e incremental. \par

\section{Desarrollo ágil}

A medida que nos adentramos en la actualidad, podemos observar como el uso de la metodología agil \cite{Hoda2018}

\section{De desarrollo ágil, a DevOps}

Para llevar a cabo cualquier tipo de proyecto, deben haber como mínimo dos partes implicadas.
El cliente y la empresa que va a llevar a cabo la vision del cliente.\par
Estas dos partes pueden ser tan complejas o simples como se quiera, pero deben existir.
Desde un proyecto llevado a cabo únicamente por una persona a un proyecto distribuido separado entre distintas zonas horarias.

La mayor de las responsabilidades del \textbf{cliente} es especificar claramente los requisitos del producto o idea que tiene en mente.

La responsabilidad de \textbf{la empresa} es de desarrollar el producto que tiene el ciente en mente lo mas fielmente posible.

La clave del éxito del desarrollo ágil, basado en iteraciones incrementales, es que establece canales de comunicación directos, estrechos e iterativos entre el cliente y la empresa.
(Existen modelos de desarrollo ágil, en el que se recomienda que el cliente comparta habitación con la empresa que esta desarrollando su producto)

Aunque, como ya se ha demostrado en empresas no directamente relacionadas con el mundo del Software como Canon, Honda, o NEC, entre otras, se puede aplicar el desarrollo ágil a cualquier sector \cite{Hirotaka1986}.
En el mundo del software, podemos hacer uso de herramientas automatizadas y frameworks para llevar a cabo los requisitos del desarrollo ágil.
%TODO: Falta una definición formal del significado de desarrollo ágil antes de llegar aquí.
De esta forma, se logra crear un proyecto que se mantiene actualizado entre las especificaciones y el desarrollo.


\subsection{Historia Seguridad en el mundo empresarial}
\subsection{Historia Introducción de la seguridad en la metodología ágil}
\subsection{Historia Inicio del fenómeno DevSecOps}

\section{presente como se gestiona la seguridad en el mundo DevSecOps}

\section{Empresas actuales, casos de éxito}

\begin{itemize}
  \item{IBM}
  \item{Netflix}
  \item{Amazon}
\end{itemize}

\subsection{Presente Aplicaciones Off the Shelf}

\begin{itemize}
  \item{Larch Darkly}
\end{itemize}


%----------
%	Bibliography
%----------	

\clearpage
\addcontentsline{toc}{chapter}{Bibliografía}

\printbibliography


%----------
%	Appendix
%----------	

% If your work includes Appendix, you can uncomment the following lines
%\chapter* {Appendix x}
%\pagenumbering{gobble} % Appendix pages are not numbered

\end{document}
